\chapter{Acknowledgements}
\label{cha:Acknowledgements}

Like many analytically minded people, I have always been drawn to the theory
and mathematics of systems. But a thesis is built on more than mathematics, and
I owe a great deal to the people who carried me through this journey. For a
long time I had a lingering curiosity about the foundations of deep learning,
though it remained vague and unformed. I am profoundly grateful to my advisor,
David, for taking on the role of guiding me since those formative days, which
helped turn an abstract interest into the concrete reality of this thesis. I am
equally grateful to my supervisor, Alexander, for helping me find my footing on
this topic. I'm very grateful to both of them for their continued guidance and
unwavering support throughout, and for anchoring me whenever I started hiding
behind complex abstractions, gently pulling me back to earth and challenging me
to find the true clarity beneath the technical language. I would also like to
thank Francesca, whose work first showed me that this topic could be both
mathematically beautiful and genuinely alive. I am deeply grateful to Prof. Dr.
Jan van Neerven for teaching me functional analysis; many of the tools I used
in this thesis began there, in a course that changed the way I saw mathematics.
I would also like to thank Prof. Dr. Frank Redig for teaching me
high-dimensional probability, a subject that ended up playing a much larger
role in the thesis than I first expected. I am also grateful to Prof. Dr.
Antonis Papapantoleon for serving on my committee. I thank my grandmother and
my parents, who have always believed in me with a steadiness I have often
borrowed when my own confidence was less reliable. I thank my girlfriend,
Radha, who was there through the long stretches when this thesis was all I
could talk about, and stayed anyway. Lastly, I also want to thank my friends, who
listened patiently while I tried to explain my thesis, sometimes with very
little warning and often with the misplaced optimism that this time it would
only take five minutes. Their love, humour, patience, and reminders that there
is life outside kernels, Fourier modes, and LaTeX error messages made this work
feel less like a solitary climb and more like something I was lucky enough to
carry among people who cared.

\vskip1cm
\begin{flushright}
	\theauthor\\
	Delft, the Netherlands \\
	\today\\
\end{flushright}

